\inhaltkopf{INHALT DES ERSTEN BANDES.}
\inhaltzeile{Einleitung}{1}
\inhaltbuch{Erstes Buch.}{Die Grundlagen.}
\inhaltabschnitt{Erster Abschnitt.}{Rationale Functionen.}
\inhaltpar{1}{Ganze Functionen}{23}
\inhaltpar{2}{Ein Satz von \gesperrt{Gauss}}{25}
\inhaltpar{3}{Division}{27}
\inhaltpar{4}{Theilung durch eine lineare Function}{30}
\inhaltpar{5}{Gebrochene Functionen; Theilbarkeit}{32}
\inhaltpar{6}{Grösster gemeinschaftlicher Theiler}{34}
\inhaltpar{7}{Producte linearer Factoren}{39}
\inhaltpar{8}{Der binomische Lehrsatz}{42}
\inhaltpar{9}{Interpolation}{44}
\inhaltpar{10}{Lösung des Interpolationsproblems durch die Differenzen}{46}
\inhaltpar{11}{Arithmetische Reihen höherer Ordnung}{47}
\inhaltpar{12}{Der polynomische Lehrsatz}{50}
\inhaltpar{13}{Derivirte Functionen}{51}
\inhaltpar{14}{Derivirte eines Productes}{54}
\inhaltpar{15}{Ganze Functionen mehrerer Veränderlicher: Formen}{56}
\inhaltpar{16}{Die Derivirten von Functionen mehrerer Variablen}{59}
\inhaltpar{17}{Das \gesperrt{Euler}'sche Theorem über homogene Functionen}{62}
\inhaltabschnitt{Zweiter Abschnitt.}{Determinanten.}
\inhaltpar{18}{Permutationen von $n$ Elementen}{64}
\inhaltpar{19}{Permutationen erster und zweiter Art}{65}
\inhaltpar{20}{Determinanten}{68}
\inhaltpar{21}{Hauptsätze über Determinanten}{70}
